\documentclass[11pt,a4paper]{article}
\usepackage{fontspec,amsmath,amssymb,amsthm,microtype}
\setmainfont{lmroman10-regular.otf}[BoldFont=lmroman10-bold.otf,ItalicFont=lmroman10-italic.otf,BoldItalicFont=lmroman10-bolditalic.otf]
\usepackage{xcolor}
\usepackage[most]{tcolorbox}
\usepackage{mathrsfs}
\providecommand{\lt}{<}
\providecommand{\gt}{>}
\setmonofont{DejaVuSansMono.ttf}[BoldFont=DejaVuSansMono-Bold.ttf,ItalicFont=DejaVuSansMono-Oblique.ttf,BoldItalicFont=DejaVuSansMono-BoldOblique.ttf,Scale=0.86]
\newfontfamily\leanfallback{STIXTwoMath-Regular.otf}[Scale=0.86]
\definecolor{leanrule}{gray}{0.62}
\definecolor{leanback}{gray}{0.962}
\definecolor{leantitle}{gray}{0.22}
\definecolor{leanmuted}{gray}{0.38}
\newtcolorbox{leanbox}[1]{enhanced,breakable,colback=leanback,colframe=leanrule,
  boxrule=0.4pt,arc=0.6mm,left=2.4mm,right=2.4mm,top=1.8mm,bottom=1.8mm,
  before skip=0.8em,after skip=1em,fonttitle=\small\bfseries,coltitle=leantitle,
  colbacktitle=leanback,titlerule=0.3pt,title={#1}}
\newlength{\leancharwidth}
\newenvironment{leancode}{\par\smallskip\begingroup\ttfamily\footnotesize\raggedright
  \settowidth{\leancharwidth}{M}\setlength{\parskip}{0pt}\setlength{\parindent}{0pt}
  \frenchspacing}{\par\endgroup\smallskip}
\newcommand{\leanline}[2]{\hangindent=\dimexpr#1\leancharwidth+4\leancharwidth\relax
  \hangafter=1\noindent\hspace*{#1\leancharwidth}#2\par}
\newcommand{\leansource}[1]{\par{\footnotesize\color{leanmuted}#1}\par}
\newcommand{\leanlabel}[1]{\par\smallskip\noindent{\small\bfseries #1}\par}
\usepackage[margin=25mm]{geometry}
\usepackage[unicode,hidelinks]{hyperref}
\hypersetup{pdfauthor={LangevinLean Project}}
\setlength{\emergencystretch}{3em}
\setlength{\parskip}{0.4em}
\setlength{\parindent}{0pt}
\allowdisplaybreaks[2]
\raggedbottom
\title{\texorpdfstring{Sharp parameter concentration and the IPLA RMS bound}{Sharp parameter concentration and the IPLA RMS bound}}
\author{LangevinLean Project}
\date{}
\begin{document}\maketitle
Draft mathematical exposition; these notes may contain mistakes.


Apply the continuous-potential concentration estimate to the parameter marginal, whose curvature has factor \(N\) from particle count. For a differentiable \(m\)-strongly convex pair potential \(U\) with \(m>0\), Prékopa–Leindler gives the marginal curvature. Its \(N\)-power target has curvature \(Nm\), so its parameter second moment is at most \(d_\theta/(Nm)\). This improves the equilibrium term in the earlier estimate while keeping the transport and discretization terms unchanged.



Earlier results used here: \cite[Section~3.19]{LangevinLean2026IPLA}.

\section{\texorpdfstring{Sharp concentration of the parameter marginal}{Sharp concentration of the parameter marginal}}


For continuous \(a\)-strongly convex potential \(f\) with minimizer \(x_\star\), the dilation argument gives

\[
\int\|x-x_\star\|^2\,\pi(dx)\le\frac d a,
\qquad\pi(dx)\propto e^{-f(x)}dx.
\]

Only continuity, strong convexity, and a minimizing point are needed; no nonsmooth integration-by-parts formula is assumed. Prékopa–Leindler gives strong convexity of the parameter marginal potential. Raising the marginal likelihood to its \(N\)-th power multiplies that curvature by \(N\). Thus

\[
W_2(\Pi_{\lambda,N}^{\theta},\delta_{\theta_{\lambda,\star}})
\le\sqrt{\frac{d_\theta}{Nm}}.
\]

This is the factor in the source’s equilibrium estimate. The earlier radius \(\sqrt{2d_\theta/(Nm)}\) was a valid weaker bound.



The parameter density is proportional to \(M(\theta)^N\) and has negative log potential \(NF(\theta)\). Its curvature is \(Nm\), with minimizer \(\theta_*\). The sharp theorem in dimension \(d_\theta\) gives \(\int\|\theta-\theta_*\|^2d\Pi_N\le d_\theta/(Nm)\). In the equivalent notation, density \(e^{-N\Phi}\) has curvature \(Nm\), so dimension \(d_\theta\) gives second moment \(d_\theta/(Nm)\). The exact projection identity transfers that same moment to the full Gibbs law; its nonnegative square root is the equilibrium root mean square radius.

\begin{leanbox}{Lean: \texttt{sharp\_\allowbreak{}marginalTarget\_\allowbreak{}secondMoment\_\allowbreak{}le}, \texttt{sharp\_\allowbreak{}gibbsTarget\_\allowbreak{}parameter\_\allowbreak{}secondMoment\_\allowbreak{}le}, \texttt{sharp\_\allowbreak{}gibbsTarget\_\allowbreak{}parameter\_\allowbreak{}concentration}}
\leansource{\texttt{sharp\_\allowbreak{}marginalTarget\_\allowbreak{}secondMoment\_\allowbreak{}le} \quad ProximalLangevin/\allowbreak{}SharpIPLA.lean, lines 23–37}
\begin{leancode}
\leanline{0}{theorem\ sharp\_marginalTarget\_secondMoment\_le\ \{N\ dθ\ dx\ :\ ℕ\}\ (hN\ :\ 0\ <\ N)}
\leanline{4}{(U\ :\ IPLA.Joint\ dθ\ dx\ →\ ℝ)\ (m\ :\ ℝ)\ (hm\ :\ 0\ <\ m)}
\leanline{4}{(hdiff\ :\ Differentiable\ ℝ\ U)\ (hsc\ :\ StrongConvexOn\ univ\ m\ U)}
\leanline{4}{(θStar\ :\ EuclideanSpace\ ℝ\ (Fin\ dθ))\ (hmax\ :\ IPLA.IsMarginalMaximizer\ U\ θStar)\ :}
\leanline{4}{(∫\ θ,\ ‖θ\ -\ θStar‖\ \textasciicircum{}\ 2\ ∂IPLA.marginalTarget\ N\ U)\ ≤}
\leanline{6}{(dθ\ :\ ℝ)\ /\ ((N\ :\ ℝ)\ *\ m)}
\end{leancode}
\leansource{\texttt{sharp\_\allowbreak{}gibbsTarget\_\allowbreak{}parameter\_\allowbreak{}secondMoment\_\allowbreak{}le} \quad ProximalLangevin/\allowbreak{}SharpIPLA.lean, lines 41–50}
\begin{leancode}
\leanline{0}{theorem\ sharp\_gibbsTarget\_parameter\_secondMoment\_le\ \{N\ dθ\ dx\ :\ ℕ\}\ (hN\ :\ 0\ <\ N)}
\leanline{4}{(U\ :\ IPLA.Joint\ dθ\ dx\ →\ ℝ)\ (m\ :\ ℝ)\ (hm\ :\ 0\ <\ m)}
\leanline{4}{(hdiff\ :\ Differentiable\ ℝ\ U)\ (hsc\ :\ StrongConvexOn\ univ\ m\ U)}
\leanline{4}{(θStar\ :\ EuclideanSpace\ ℝ\ (Fin\ dθ))\ (hmax\ :\ IPLA.IsMarginalMaximizer\ U\ θStar)\ :}
\leanline{4}{(∫\ z\ :\ IPLA.State\ N\ dθ\ dx,\ ‖IPLA.parameterCLM\ z\ -\ θStar‖\ \textasciicircum{}\ 2}
\leanline{6}{∂IPLA.gibbsTarget\ U)\ ≤\ (dθ\ :\ ℝ)\ /\ ((N\ :\ ℝ)\ *\ m)}
\end{leancode}
\leansource{\texttt{sharp\_\allowbreak{}gibbsTarget\_\allowbreak{}parameter\_\allowbreak{}concentration} \quad ProximalLangevin/\allowbreak{}SharpIPLA.lean, lines 53–71}
\begin{leancode}
\leanline{0}{theorem\ sharp\_gibbsTarget\_parameter\_concentration\ \{N\ dθ\ dx\ :\ ℕ\}\ (hN\ :\ 0\ <\ N)}
\leanline{4}{(U\ :\ IPLA.Joint\ dθ\ dx\ →\ ℝ)\ (m\ :\ ℝ)\ (hm\ :\ 0\ <\ m)}
\leanline{4}{(hdiff\ :\ Differentiable\ ℝ\ U)\ (hsc\ :\ StrongConvexOn\ univ\ m\ U)}
\leanline{4}{(θStar\ :\ EuclideanSpace\ ℝ\ (Fin\ dθ))\ (hmax\ :\ IPLA.IsMarginalMaximizer\ U\ θStar)\ :}
\leanline{4}{IPLA.projectedRMS\ (@IPLA.gibbsTarget\ N\ dθ\ dx\ U)\ IPLA.parameterCLM\ θStar\ ≤}
\leanline{6}{ENNReal.ofReal\ (Real.sqrt\ ((dθ\ :\ ℝ)\ /\ ((N\ :\ ℝ)\ *\ m)))}
\end{leancode}
\leanlabel{How the Lean proof follows the argument}
Formal correspondence: sharp\_marginalTarget\_secondMoment\_le : Sharp equilibrium second moment of the actual powered marginal density. sharp\_gibbsTarget\_parameter\_secondMoment\_le : The exact parameter-marginal identity transfers sharp concentration to the full lifted Gibbs measure. sharp\_gibbsTarget\_parameter\_concentration : The parameter RMS radius with the sharp dimension constant.



\texttt{sharp\_\allowbreak{}marginal\allowbreak{}Target\_\allowbreak{}second\allowbreak{}Moment\_\allowbreak{}le}: the intermediate facts include \texttt{hmode}, \texttt{h}, in source order; the source steps apply or rewrite \texttt{IPLA.\allowbreak{}marginal\allowbreak{}Target\_\allowbreak{}eq\_\allowbreak{}gibbs\allowbreak{}Measure}, \texttt{finrank\_\allowbreak{}euclidean\allowbreak{}Space}, \texttt{Fintype.\allowbreak{}card\_\allowbreak{}fin}.



\texttt{sharp\_\allowbreak{}gibbs\allowbreak{}Target\_\allowbreak{}parameter\_\allowbreak{}second\allowbreak{}Moment\_\allowbreak{}le}: the intermediate fact is \texttt{h}, in source order; the source steps apply or rewrite \texttt{IPLA.\allowbreak{}gibbs\allowbreak{}Target\_\allowbreak{}parameter\_\allowbreak{}marginal}.



\texttt{sharp\_\allowbreak{}gibbs\allowbreak{}Target\_\allowbreak{}parameter\_\allowbreak{}concentration}: the intermediate fact is \texttt{hi}, in source order; the source steps apply or rewrite \texttt{IPLA.\allowbreak{}projected\allowbreak{}RMS\_\allowbreak{}le\_\allowbreak{}of\_\allowbreak{}second\allowbreak{}Moment}, \texttt{Real.\allowbreak{}sqrt\_\allowbreak{}nonneg}, \texttt{Real.\allowbreak{}sq\_\allowbreak{}sqrt}, \texttt{sharp\_\allowbreak{}gibbs\allowbreak{}Target\_\allowbreak{}parameter\_\allowbreak{}second\allowbreak{}Moment\_\allowbreak{}le}.
\end{leanbox}


Earlier results used here: \cite[Section~3.26.13]{LangevinLean2026ULASC}; \cite[Section~3.26.15]{LangevinLean2026ULASC}; \cite[Section~3.26.3]{LangevinLean2026ULASC}; \cite[Section~3.1.1]{LangevinLean2026IPLA}; \cite[Section~3.1.3]{LangevinLean2026IPLA}; \cite[Section~3.1.6]{LangevinLean2026IPLA}; \cite[Section~3.12.1]{LangevinLean2026IPLA}; \cite[Section~3.13.1]{LangevinLean2026IPLA}; \cite[Section~3.3.1]{LangevinLean2026IPLA}; \cite[Section~3.3.2]{LangevinLean2026IPLA}; \cite[Section~3.9.9]{LangevinLean2026IPLA}.

\section{\texorpdfstring{The Gaussian particle estimate with sharp equilibrium radius}{The Gaussian particle estimate with sharp equilibrium radius}}


Assume \(N>0\), \(\gamma>0\), \(\gamma\le2/(m+L)\), independent Gaussian innovations, and \(L^2\) initialization. The scaled Gaussian Euler estimate, parameter projection, and the sharp equilibrium moment give

\[\operatorname{RMS}_n\le\sqrt{\frac{d_\theta}{Nm}}+(1-m\gamma)^nW_2(\rho_0,\Pi_N)+\frac{33}{20}\frac Lm\sqrt{d_x+d_\theta/N}\sqrt\gamma.\]

For any coupling of the current state and Gibbs target, projection does not increase quadratic cost. Minkowski adds the Gibbs parameter error about the center. Taking the coupling infimum gives transport plus the sharp radius, and substitution yields the displayed estimate. If the parameter dimension is zero its error is zero; if the state is a singleton its transport is zero. Thus the same inequality holds for arbitrary nonnegative parameter and latent dimensions.

\begin{leanbox}{Lean: \texttt{sharp\_\allowbreak{}ipla\_\allowbreak{}parameter\_\allowbreak{}rms\_\allowbreak{}nonempty}, \texttt{sharp\_\allowbreak{}ipla\_\allowbreak{}parameter\_\allowbreak{}rms}}
\leansource{\texttt{sharp\_\allowbreak{}ipla\_\allowbreak{}parameter\_\allowbreak{}rms\_\allowbreak{}nonempty} \quad ProximalLangevin/\allowbreak{}SharpIPLA.lean, lines 73–124}
\begin{leancode}
\leanline{0}{private\ theorem\ sharp\_ipla\_parameter\_rms\_nonempty}
\leanline{4}{\{Ω\ :\ Type*\}\ [MeasurableSpace\ Ω]\ (Q\ :\ Measure\ Ω)\ [IsProbabilityMeasure\ Q]}
\leanline{4}{(N\ dθ\ dx\ :\ ℕ)\ [Nonempty\ (Fin\ dθ\ ⊕\ (Fin\ N\ ×\ Fin\ dx))]\ (hN\ :\ 0\ <\ N)}
\leanline{4}{(U\ :\ IPLA.Joint\ dθ\ dx\ →\ ℝ)\ (m\ γ\ :\ ℝ)\ (L\ :\ ℝ≥0)}
\leanline{4}{(hdiff\ :\ Differentiable\ ℝ\ U)\ (hsc\ :\ StrongConvexOn\ univ\ m\ U)}
\leanline{4}{(hLip\ :\ LipschitzWith\ L\ (gradient\ U))\ (hm\ :\ 0\ <\ m)}
\leanline{4}{(hγ\ :\ 0\ <\ γ)\ (hsmall\ :\ γ\ ≤\ 2\ /\ (m\ +\ (L\ :\ ℝ)))}
\leanline{4}{(θStar\ :\ EuclideanSpace\ ℝ\ (Fin\ dθ))\ (hmax\ :\ IPLA.IsMarginalMaximizer\ U\ θStar)}
\leanline{4}{(initial\ :\ Lp\ (IPLA.State\ N\ dθ\ dx)\ 2\ Q)}
\leanline{4}{(noise\ :\ ℕ\ →\ Lp\ (IPLA.State\ N\ dθ\ dx)\ 2\ Q)}
\leanline{4}{(hGaussian\ :\ IPLA.GaussianInnovations\ Q\ initial\ noise)\ (n\ :\ ℕ)\ :}
\leanline{4}{Real.sqrt\ (∫\ z,\ ‖IPLA.parameterCLM\ z\ -\ θStar‖\ \textasciicircum{}\ 2}
\leanline{6}{∂IPLA.laws\ Q\ U\ γ\ L\ hdiff\ hLip\ initial\ noise\ n)\ ≤}
\leanline{6}{Real.sqrt\ ((dθ\ :\ ℝ)\ /\ ((N\ :\ ℝ)\ *\ m))\ +}
\leanline{6}{(1\ -\ m\ *\ γ)\ \textasciicircum{}\ n\ *\ ULA.wasserstein2Real\ (Q.map\ initial)\ (IPLA.gibbsTarget\ U)\ +}
\leanline{6}{(33\ /\ 20\ :\ ℝ)\ *\ ((L\ :\ ℝ)\ /\ m)\ *}
\leanline{8}{Real.sqrt\ ((dx\ :\ ℝ)\ +\ (dθ\ :\ ℝ)\ /\ (N\ :\ ℝ))\ *\ Real.sqrt\ γ}
\end{leancode}
\leansource{\texttt{sharp\_\allowbreak{}ipla\_\allowbreak{}parameter\_\allowbreak{}rms} \quad ProximalLangevin/\allowbreak{}SharpIPLA.lean, lines 128–158}
\begin{leancode}
\leanline{0}{theorem\ sharp\_ipla\_parameter\_rms}
\leanline{4}{\{Ω\ :\ Type*\}\ [MeasurableSpace\ Ω]\ (Q\ :\ Measure\ Ω)\ [IsProbabilityMeasure\ Q]}
\leanline{4}{(N\ dθ\ dx\ :\ ℕ)\ (hN\ :\ 0\ <\ N)\ (U\ :\ IPLA.Joint\ dθ\ dx\ →\ ℝ)}
\leanline{4}{(m\ γ\ :\ ℝ)\ (L\ :\ ℝ≥0)\ (hdiff\ :\ Differentiable\ ℝ\ U)}
\leanline{4}{(hsc\ :\ StrongConvexOn\ univ\ m\ U)\ (hLip\ :\ LipschitzWith\ L\ (gradient\ U))}
\leanline{4}{(hm\ :\ 0\ <\ m)\ (hγ\ :\ 0\ <\ γ)\ (hsmall\ :\ γ\ ≤\ 2\ /\ (m\ +\ (L\ :\ ℝ)))}
\leanline{4}{(θStar\ :\ EuclideanSpace\ ℝ\ (Fin\ dθ))\ (hmax\ :\ IPLA.IsMarginalMaximizer\ U\ θStar)}
\leanline{4}{(initial\ :\ Lp\ (IPLA.State\ N\ dθ\ dx)\ 2\ Q)}
\leanline{4}{(noise\ :\ ℕ\ →\ Lp\ (IPLA.State\ N\ dθ\ dx)\ 2\ Q)}
\leanline{4}{(hGaussian\ :\ IPLA.GaussianInnovations\ Q\ initial\ noise)\ (n\ :\ ℕ)\ :}
\leanline{4}{Real.sqrt\ (∫\ z,\ ‖IPLA.parameterCLM\ z\ -\ θStar‖\ \textasciicircum{}\ 2}
\leanline{6}{∂IPLA.laws\ Q\ U\ γ\ L\ hdiff\ hLip\ initial\ noise\ n)\ ≤}
\leanline{6}{Real.sqrt\ ((dθ\ :\ ℝ)\ /\ ((N\ :\ ℝ)\ *\ m))\ +}
\leanline{6}{(1\ -\ m\ *\ γ)\ \textasciicircum{}\ n\ *\ ULA.wasserstein2Real\ (Q.map\ initial)\ (IPLA.gibbsTarget\ U)\ +}
\leanline{6}{(33\ /\ 20\ :\ ℝ)\ *\ ((L\ :\ ℝ)\ /\ m)\ *}
\leanline{8}{Real.sqrt\ ((dx\ :\ ℝ)\ +\ (dθ\ :\ ℝ)\ /\ (N\ :\ ℝ))\ *\ Real.sqrt\ γ}
\end{leancode}
\leanlabel{How the Lean proof follows the argument}
For \(D=d_\theta+Nd_x\), the lifted state has score constant \(NL\), convexity \(Nm\) and step \(\gamma/N\). The shared Gaussian ULA theorem controls its \(W_2\) error, while the parameter projection has norm at most one. The projected RMS therefore is at most the sharp equilibrium scale \(\sqrt{d_\theta/(Nm)}\) plus that transport error. The nonempty-dimensional case proves this combination directly; the final declaration splits off zero-dimensional cases so that no unspoken positive-dimension premise remains.



Formal correspondence: sharp\_ipla\_parameter\_rms : A complete actual Gaussian IPLA parameter bound with sharp concentration. This theorem includes zero parameter/\allowbreak{}latent dimensions.



\texttt{sharp\_\allowbreak{}ipla\_\allowbreak{}parameter\_\allowbreak{}rms\_\allowbreak{}nonempty}: the intermediate facts include \texttt{hn}, \texttt{hfinite}, \texttt{hs}, \texttt{hp}, \texttt{hmoment}, in source order; the source steps apply or rewrite \texttt{IPLA.\allowbreak{}scaled\_\allowbreak{}step\_\allowbreak{}iff}, \texttt{IPLA.\allowbreak{}gibbs\allowbreak{}Target\_\allowbreak{}eq\_\allowbreak{}gibbs\allowbreak{}Measure}, \texttt{IPLA.\allowbreak{}scaled\_\allowbreak{}bias\_\allowbreak{}eq}, \texttt{h\ensuremath{\gamma}.\allowbreak{}le}, \texttt{IPLA.\allowbreak{}projected\allowbreak{}RMSReal\_\allowbreak{}eq\_\allowbreak{}sqrt\_\allowbreak{}integral}. The remaining scalar rearrangements are checked by \texttt{field\_\allowbreak{}simp}, \texttt{linarith}.



\texttt{sharp\_\allowbreak{}ipla\_\allowbreak{}parameter\_\allowbreak{}rms}: the cases are separated by \texttt{by\_\allowbreak{}cases hi : Nonempty (Fin d\ensuremath{\theta} ⊕ (Fin N \ensuremath{\times} Fin dx))}; the intermediate facts include \texttt{h\ensuremath{\theta}}, \texttt{hx}, in source order; the source steps apply or rewrite \texttt{sharp\_\allowbreak{}ipla\_\allowbreak{}parameter\_\allowbreak{}rms\_\allowbreak{}nonempty}, \texttt{Nat.\allowbreak{}cast\_\allowbreak{}zero}, \texttt{mul\_\allowbreak{}zero}.
\end{leanbox}


Earlier results used here: \cite[Section~3.17.4]{LangevinLean2026ULASC}; \cite[Section~3.26.3]{LangevinLean2026ULASC}; \cite[Section~3.45.6]{LangevinLean2026ULASC}; \cite[Section~3.50.2]{LangevinLean2026ULASC}; \cite[Section~3.50.3]{LangevinLean2026ULASC}; \cite[Section~3.1.1]{LangevinLean2026IPLA}; \cite[Section~3.1.3]{LangevinLean2026IPLA}; \cite[Section~3.1.4]{LangevinLean2026IPLA}; \cite[Section~3.1.6]{LangevinLean2026IPLA}; \cite[Section~3.16.3]{LangevinLean2026IPLA}; \cite[Section~3.18.2]{LangevinLean2026IPLA}; \cite[Section~3.3.1]{LangevinLean2026IPLA}; \cite[Section~3.3.2]{LangevinLean2026IPLA}; \cite[Section~3.3.3]{LangevinLean2026IPLA}; \cite[Section~3.5.1]{LangevinLean2026IPLA}; \cite[Section~3.7.1]{LangevinLean2026IPLA}.

\begin{thebibliography}{ULA-LSI}
\bibitem[ULA-SC]{LangevinLean2026ULASC}
LangevinLean Project. \emph{Unadjusted Langevin under Strong Convexity}. LL-01, 2026. Draft report, content revision d34ead06987c. \url{https://odakyildiz.github.io/LangevinLean/docs/pdf/UnadjustedLangevin/FullProof.pdf}.
\bibitem[IPLA-SC]{LangevinLean2026IPLA}
LangevinLean Project. \emph{Interacting Particle Langevin under Strong Convexity}. LL-03, 2026. Draft report, content revision ec0201df2207. \url{https://odakyildiz.github.io/LangevinLean/docs/pdf/IPLA/FullProof.pdf}.
\end{thebibliography}

\end{document}
